\newcommand{\dimK}{\dim_{\text{Krull}}}
\newcommand{\dimH}{\dim_{\text{Hilbert}}}
\newcommand{\dimP}{\dim_{\text{primary}}}
\newcommand{\height}{\operatorname{ht}}
\chapter{Dimension Theory}
\label{chap:dimension}
We have met the notion of dimension of a ring in
\Cref{sec:dim_intro}.
To reproduce the definition, recall that
for any ring $A$, the \vocab{Krull dimension}
\notation{$\dimK A$}{Krull dimension}
$\dimK A$ is the maximum nonnegative integer
$n$ such that there exists a chain of prime ideals
$$\pp_0\subsetneq \pp_1 \subsetneq \pp_2
\subsetneq \dots\subsetneq \pp_n \ \subset\ A.$$
This notion of dimension is notoriously difficult
to deduce its properties.
For example, proving that $\dim \CC[x_1,\dots,x_n]=n$,
which we did in \Cref{sec:dim_intro}, was a nontrivial task.

To give a more robust theory of dimension,
we first note that we can reduce to the study of local ring.
Since any chain of prime ideals should terminate at a maximal ideal,
we deduce that
$$\dimK A = \max_{\mm\text{ maximal}} \dimK A_\mm 
= \max_{\pp\text{ prime}} \dimK A_\pp.$$
Thus, the study of $\dimK A$ is reduced to the study
of $\dimK A$ of a local ring.

Given a Noetherian local ring $A$ with maximal ideal $\mm$,
we will define two other types of dimensions:
\begin{itemize}
    \item the \vocab{primary ideal dimension}
    $\dimP A$,
    which is the minimum number of generators
    across all $\mm$-primary ideals;
    \item the \vocab{Hilbert dimension}
    $\dimH A$, which is more complicated and will be discussed
    in \Cref{sec:hilbert_dim}.
\end{itemize}
Then, the main theorem is 
$\dimK A = \dimH A = \dimP A$ for any local ring $A$.

Before discussing dimension in more detail,
let us close with the following warning.
\begin{warn}
Noetherian rings need not have finite dimension.
In addition, rings with finite dimension need 
not be Noetherian.
\Cref{prob:dim_pathology} provides counterexamples.
\end{warn}
\section{Preliminary: Length of a Module}
\begin{definition}[Composition Series]
    Given a ring $A$ and $A$-module $M$,
    a \vocab{composition series} of a module $M$
    is a chain
    $$M=M_0\supsetneq M_1\supsetneq M_2
    \supsetneq \dots\supsetneq M_n=0$$
    such that $M_k/M_{k+1}$ is a \vocab{simple module}
    (i.e., have no submodule other than $0$
    and $M_k/M_{k+1}$).
\end{definition}
\begin{exercise}
    Prove that $M$ is simple if and only if 
    $M=A/\mm$ for a maximal ideal $\mm$.
\end{exercise}
This definition is an analogue of vector space.
In particular, if $A=k$ is a field and $M\simeq k^n$,
then we have the composition series
$$k^d \supset k^{d-1}\supset\dots\supset k^1\supset k^0=\{0\}.$$
However, in contrast to vector spaces,
composition series of a module does not always exist.
For example, if $M=A=\mathbb Z$,
then anything in the chain must be isomorphic to $\mathbb Z$,
so we cannot get it down to zero.
\begin{definition}[Length]
    For any module $M$, define the
    \vocab{length} $l(M)$ as the minimum length
    of a composition of a composition series of $M$.
    (let $l(M)=\infty$ if no finite composition series exists).
    \notation{$l(M)$}{length of a module}
\end{definition}
For example, we have $l(\mathbb C^d) = d$
(as a $\CC$-module) and $l(\ZZ)=\infty$ 
(as a $\ZZ$-module).

It is not obvious that length is well-defined.
Thus, we prove it.
\begin{proposition}[Length is well-defined]
    \label{prop:length_is_well_defined}
    If $l(M)=n$, then
    every composition series has length $n$,
    and any chain of submodules can be extend to
    a composition series.
\end{proposition}
To prove this, we need the following lemma.
\begin{lemma}[Length of a proper submodule]
    If $N\subsetneq M$
    and $l(M)<\infty$, then $l(N) < l(M)$.
\end{lemma}
\begin{proof}
    Let $n=l(M)$.
    Take a composition series
    $$M=M_0\supsetneq M_1\supsetneq \dots\supsetneq M_n =0.$$
    Let $N_i = M_i\cap N$.
    Then, we have
    $$N_0\supset N_1\supset\dots\supset N_n.$$
    Then, we can see that $N_k/N_{k+1} \subset M_k/M_{k+1}$,
    so either $N_k/N_{k+1}=0$ or $N_k/N_{k+1}$ is simple.
    Thus, $\{N_k\}$ is a composition series after
    erasing equal module.
    Therefore, we already have that $l(N)\leq l(M)$.

    If $l(N)=l(M)$, then the only possibility is that
    we did not do the erasing in the paragraph.
    In particular, we have that for each $k$,
    $N_k/N_{k+1} = M_k/M_{k+1}$.
    Thus, we have
    $$\begin{array}{rcl}
	N_{n-1}/0  = M_{n-1}/0 &\implies & N_{n-1}=M_{n-1}. \\[4pt]
	N_{n-2}/N_{n-1} = M_{n-2}/M_{n-1}
		&\implies & N_{n-2}=M_{n-2} \\[4pt]
		& \vdots & \\[4pt]
	      &\implies & N=M,
    \end{array}$$
    which is a contradiction to $N\subsetneq M$.
\end{proof}
\begin{proof}[Proof of \Cref{prop:length_is_well_defined}]
    Given any chain of length $k$, we have
    $$M = M_0\supsetneq M_1\supsetneq M_2\supsetneq
    \dots\supsetneq M_k = 0,$$
    then by the lemma, we have
    $$l(M_0) > l(M_1) > \dots > l(M_k)\geq 0,$$
    so $l(M)\geq k$.
    In particular, any composition series has length $\geq l(M)$.
    Moreover, any chain can be extended to a composition series.
\end{proof}
We now prove some more properties of a length.
\begin{proposition}[Length behaves well under short exact sequence]
    \label{prop:length_exact}
    Suppose we have a short exact sequence
    $0\to N\to M\to P\to 0$,
    then $l(M) = l(N)+l(P)$.
\end{proposition}
\begin{proof}
    By \Cref{prop:length_is_well_defined},
    given a composition series of $N$ and $P$,
    it suffices to find a composition of $M$.
    To that end, suppose we have composition series
    \begin{align*}
        N = N_0\supsetneq N_1\supsetneq \dots \supsetneq N_k=0 \\
	P = P_0\supset P_1\supset\dots\supset P_\ell = 0.
    \end{align*}
    Name the map $\pi : N\to P$.
    Then, we have a composition series of length $k+\ell$:
    \[0\subset N_k\subset N_{k-1}\subset
    \dots\subset N_0=N=\pi^{-1}(0)
    \subset \pi^{-1}P_{\ell-1}
	\subset \pi^{-1}P_{\ell-2}
    \subset\dots\subset \pi^{-1}P_0 = P.\qedhere\]
\end{proof}
\begin{corollary}
    \label{cor:length_exact}
    If there is a long exact sequence
    $$0\longto M_1 \longto M_2 \longto \dots\longto M_n\longto 0,$$
    then 
    $$\sum_{i=1}^n (-1)^i l(M_i) = 0.$$
\end{corollary}
\begin{proof}
    Break the long exact sequence into short exact sequences:
    Let $\phi_i$ be the map from $M_i\to M_{i+1}$.
    Then, we have the short exact sequence
    $$0\longto \Ker\phi_i\longto M_i\longto \Im\phi_i\longto 0.$$
    Thus, by \Cref{prop:length_exact}, we have
    $$l(M_i) = l(\Im\phi_i) + l(\Ker\phi_i)
    = l(\Im\phi_i) + l(\Im\phi_{i-1}).$$
    Substituting this into the sum gives the desired result.
\end{proof}
\begin{proposition}
    $M$ has finite length if and only if $M$ is both
    Noetherian and Artinian.
\end{proposition}
\begin{proof}
    $(\Rightarrow)$
    Suppose $M_1\subset M_2\subset\dots$,
    then $l(M_k)$ is increasing but bounded by $l(M)$.
    Same goes for descending chain.

    $(\Leftarrow)$
    Define
    $$\begin{array}{rcl}
	M_0 &=& M \\[4pt]
	M_1 &=& \text{maximal proper submodule of } M \\[4pt]
	M_2 &=& \text{maximal proper submodule of } M_1 \\[4pt]
	    &\vdots&
    \end{array}$$
    These exist since $M$ is Noetherian.
    We get a descending chain
    $M_0\supset M_1\supset M_2\supset\dots$,
    so since $M$ is Artinian,
    this chain eventually terminates.
\end{proof}
\begin{corollary}
    \label{cor:artinian_finite_length}
    Let $A$ be an Artinian ring,
    then any finitely-generated $A$-module
    has finite length.
\end{corollary}
%
%
%
%
\section{Hilbert's Dimension}
\label{sec:hilbert_dim}
\subsection{Definition}
Before going into the definition, let us see 
some examples first.
\begin{example}[Dimension of $k[x_1,\dots,x_n]$]
    Let $A = k[x_1,\dots,x_n]_{(x_1,\dots,x_n)}$
    (localized at maximal ideal $(x_1,\dots,x_n)$).
    Then, $A$ is a local ring with maximal ideal 
    $\mm=(x_1,\dots,x_n)$ and $A/\mm=k$.
    As a $k$-vector space, $\mm^i/\mm^{i+1}$ 
    is a linear combinations of terms $x_1^{a_1}\cdots x_n^{a_n}$
    where $a_1+\dots+a_n = i$. Thus,
    \begin{align*}\dim_k \mm^i/\mm^{i+1}
    &= \#\{(a_1,\dots,a_n) : a_1,\dots,a_n\geq 0,
    a_1+\dots+a_n=i\} \\
    &= \binom{i+n-1}{n-1},
    \end{align*}
    by the ``stars and bars trick''. 
    This is a polynomial of degree $n-1$, one less than
    the dimension.
\end{example}
\begin{example}[Artinian local ring]
    On the other hand, if $A$ is the Artinian local ring,
    which has dimension $0$, then $\dim_k \mm^i/\mm^{i+1}$
    is not a polynomial.
    Instead, we only know (by \Cref{thm:artinian_local_ring})
    that $\mm^N=0$ for some $N$, so
    $\dim_k \mm^i/\mm^{i+1}=0$ for all sufficiently large $i$.

    If we regard zero polynomial as having ``degree $-1$'',
    then again, $\dim_k \mm^i/\mm^{i+1}$
    is a polynomial of degree one less than the dimension.
\end{example}
Given a Noetherian local ring $A$ with maximal ideal $\mm$,
the above two examples lead us to consider the sequence
$$\ell_i = \dim_k \mm^i/\mm^{i+1}.$$
A sequence $(\ell_i)_{i\geq 0}$ is said to be 
\vocab{eventually polynomial} if there is a polynomial
$p\in \RR[x]$ such that $\ell_i=p(i)$ for all sufficiently 
large $i$. 
If such $p$ exists (which must then be unique), 
we say that the sequence $(\ell_i)_{i\geq 0}$
is eventually polynomial with degree $\deg p$.
\textbf{For the purpose of these notes,
the degree of zero polynomial is $-1$.}
\begin{theorem}[Hilbert dimension]
    \label{thm:hilbert_dim}
    If $A$ is a Noetherian local ring with maximal ideal $\mm$,
    then the \vocab{Hilbert-Samuel function}
    $$\ell_i = \dim_{A/\mm} \mm^i/\mm^{i+1}$$
    is eventually polynomial.
    
    If it is eventually polynomial with degree $d$,
    then we define the \vocab{Hilbert dimension}
    \notation{$\dimH A$}{Hilbert dimension}
    $$\dimH A := d+1.$$
\end{theorem}
We will prove this theorem in a greater generality
(that we will need later) in the next section.
\subsection{Proof of Polynomiality}
\paragraph{Setup} Let $B_\bullet = B_0\oplus B_1\oplus\dots$
be a graded ring such that
\begin{itemize}
    \item $B_0$ is Artinian.
    \item $B_\bullet$ is generated (as an $B_0$-algebra)
    by finitely many elements in $B_1$.
\end{itemize}
Let $M_\bullet = M_0\oplus M_1\oplus\dots$ be a finitely-generated
$B_\bullet$-module.
This implies that, for each $i$, 
$M_i$ is finitely-generated $B_0$-module,
hence has finite length (by \Cref{cor:artinian_finite_length}).

\paragraph{Main Theorem}
Now, we state the main theorem.
\begin{theorem}[Hilbert-Serre]
    \label{thm:hilbert_dim_module}
    Assuming the above setup, the sequence
    $$\ell_i := l_{B_0}(M_i)$$
    is eventually polynomial.
    (The subscript in $l_{B_0}$ means the length
    as $B_0$-module.)
    
    Furthermore, if $B_\bullet$ is generated by 
    $m$ elements in $B_1$, 
    then it is eventually polynomial of degree at most $m-1$.
\end{theorem}
To deduce \Cref{thm:hilbert_dim} from \Cref{thm:hilbert_dim_module},
we plug in the \textbf{associated graded ring}
(\Cref{def:gr})
$$B_\bullet = M_\bullet = \gr_\mm A = \bigoplus_{i\geq 0} \mm^i/\mm^{i+1},$$
and clearly the generators of $\mm$,
which sits in $\mm/\mm^2$, are sufficient
to generate $\mm^i/\mm^{i+1}$ for all $i$,
so $B_\bullet$ is finitely-generated at degree $1$.

(As a consequence, we get that the Hilbert dimension 
is at most the minimum number of generators of $\mm$.
Soon, we will prove a stronger result
replacing $\mm$ with any $\mm$-primary ideal,
which will play a key role in proving dimension theorem.)
\begin{proof}
    We first note the following elementary fact:
    if $c\geq 0$ is an integer such that
    the sequence $\ell'_i = \ell_{i+1}-\ell_i$
    is eventually polynomial with degree $d$,
    then the sequence $\ell_i$ is eventually polynomial
    with degree $d+1$.
    (This follows from ``finite difference''.)

    We use induction on $m$,
    the number of generators. The base case $m=0$ is clear
    because then, $M_\bullet$ is finitely-generated 
    $B_0$-module, so $M_n=0$ for all large $n$.
    Thus, $\ell_i$ is eventually a zero polynomial,
    which has degree $-1$.

    Now, assume that the result is known for $m-1$.
    Let $x$ be one of the generators, and define
    $$K_i = \Ker(M_i\stackrel{\cdot x}\longto M_{i+1})
    \quad\text{and}\quad 
    L_{i+1} = M_{i+1}/xM_i,$$ 
    so that we have the exact sequence
    $$0\longto K_i \longto M_i 
    \stackrel{\cdot x}\longto M_{i+1} 
    \longto L_{i+1}\longto 0,$$
    which gives (by \Cref{cor:length_exact})
    $$l_{B_0}(K_i) - l_{B_0}(M_i) + l_{B_0}(M_{i+1}) - l_{B_0}(L_i) = 0.$$
    Note that $K_\bullet$ and $L_\bullet$ are graded 
    $B_\bullet/(x)$-module. Note that $B_1/(x)$
    is generated by other $m-1$ generators,
    so by induction hypothesis,
    $l_{B_0}(K_\bullet)$ and $l_{B_0}(L_\bullet)$
    is eventually polynomial of degree at most $m-2$.
    Thus, $l_{B_0}(M_{i+1}) - l_{B_0}(M_i)
    = l_{B_0}(L_i) - l_{B_0}(K_i)$
    is eventually polynomial of degree at most $m-2$.
    The elementary fact at the beginning concludes.
\end{proof}
We note the following useful fact.
\begin{proposition}[Alternative definition of $\dimH$]
    \label{prop:alternative_length}
    If $Q$ is an ideal of $A$, then the following are equivalent:
    \begin{alphenum}
    \item $\ell_i = l_{A/Q}(Q^i/Q^{i+1})$
    is eventually polynomial of degree $d$.
    \item $p_i = l_{A/Q}(A/Q^i)$
    is eventually polynomial of degree $d+1$.
    \end{alphenum}
\end{proposition}
\begin{proof}
The exact sequence
$$0\longto Q^i/Q^{i+1} \longto A/Q^{i+1}
\longto A/Q^i\longto 0$$
gives 
$$l_{A/Q}(A/Q^{i+1}) - l_{A/Q}(A/Q^i) = l(Q^{i+1}/Q^i),$$
so $p_i-p_{i+1} = \ell_i$, 
which implies the result by finite difference.
\end{proof}
By plugging in $Q=\mm$, the Hilbert dimension 
can also be defined as the degree of $l(A/\mm^i)$.

\section{Primary Ideal Dimension and Dimension Theorem}
\subsection{Primary Ideal Dimension}
We make the third definition of dimension:
\begin{definition}[Primary ideal dimension]
    For any local ring $A$ with maximal ideal $\mm$,
    we define the \vocab{primary ideal dimension},
    denoted $\dimP A$, to be the minimum
    nonnegative integer $s$ for which there exists
    an $\mm$-primary ideal $Q$ generated by 
    $\leq s$ elements.
    \notation{$\dimP A$}{primary ideal dimension}
\end{definition}
\begin{example}[Examples of primary ideal dimension]
    \begin{itemize}
    \item Consider $A=\CC[x_1,\dots,x_n]_{(x_1,\dots,x_n)}$,
    which is a local ring with maximal ideal
    $\mm = (x_1,\dots,x_n)$. 
    The ideal $\mm$ is $\mm$-primary and has $n$ generators,
    so we get that $\dimP A \leq n$.
    \item Let $A$ be an Artinian local ring 
    with maximal ideal $\mm$.
    By \Cref{thm:artinian_local_ring}, $\mm^N=(0)$
    for some $N$, so $(0)$ is  $\mm$-primary,
    which needs $0$ generators, so $\dimP A = 0$.
    \item Consider
    $$A = \frac{\CC[x,y]_{(x,y)}}{(xy)}$$
    which is a local ring with maximal ideal $\mm=(x,y)$.
    This ideal is generated by two elements.
    However, $Q=(x+y)$ is $\mm$-primary because
    $x^2=x(x+y)\in Q$ and $y^2=y(x+y)\in Q$,
    so $\sqrt Q = \mm$, so by \Cref{prop:max_ideal_primary},
    $Q$ is $\mm$-primary.

    In particular, this implies that 
    $\dimP A \leq 1$, and it is not enough to count
    the number of generators of $\mm$.
    \end{itemize}
    Notice that $\dimP$ is easier to provide
    an \textbf{upper bound}.
    In contrast, $\dimK$ is easier to provide
    a \textbf{lower bound}.
\end{example}
There is no surprise that the main theorem
that we will prove is the following.
\begin{theorem}[Dimension Theorem]
    \label{thm:dimension}
    For any Noetherian local ring $A$, we have
    $$\dimK A = \dimH A = \dimP A.$$
    Hence, we can safely call them all $\dim A$.
\end{theorem}

Proving this will occupy the rest of this subsection.
We will show a chain of inequalities
$$\dimP A \stackrel{\text{(1)}}\geq \dimH A 
\stackrel{\text{(2)}}\geq \dimK A 
\stackrel{\text{(3)}}\geq \dimP A.$$
\Cref{subsec:dimP_dimH} will prove (1),
\Cref{subsec:dimH_dimK} will prove (2),
and \Cref{subsec:dimK_dimP} will prove (3).
Each inequality takes some amount of work.

For the rest of this section,
$A$ is a Noetherian local ring, $\mm$ is the maximal ideal,
and $k=A/\mm$ is the residue field.

\subsection{The Primary Ideal Dimension is at least the Hilbert Dimension}
\label{subsec:dimP_dimH}
We prove $\dimP A\geq \dimH A$. 
In plain terms, given a primary ideal $Q$
with $m$ generators, we have to show that 
$\dim_k \mm^{i+1}/\mm^i$ is eventually polynomial
of degree at most $m-1$.

The strategy for this is to apply 
\Cref{thm:hilbert_dim_module} to graded modules 
related to $Q$, and show that the length sequence
must be within constant factor of $\dim_k \mm^{i+1}/\mm^i$.
This will show that the degrees are equal,
and will provide a degree upper bound of
$\dim_k \mm^{i+1}/\mm^i$.



To do that, we let $M$ and $B$ be the associated graded ring
corresponding to the $Q$-adic filtration
$$M=B=\gr_Q(A)=\bigoplus_{i\geq 0} Q^i / Q^{i+1}.$$
We check the hypothesis of \Cref{thm:hilbert_dim_module}.
\begin{itemize}
    \item $B_0 = A/Q$ is Artinian because \Cref{thm:artinian_local_ring} 
    (condition (c)) applies.
    \item $B$ is clearly generated in degree $1$
    by the $m$ generators of $Q$.
\end{itemize}
Thus, by \Cref{thm:hilbert_dim_module}, the sequence
$l_{A/Q}(Q^i/Q^{i+1})$ is eventually polynomial
with degree $d\leq m-1$.

Now, let $d'$ be the degree of $\dim_k\mm^{i+1}/\mm^i$.
We claim that $d=d'$, which will finish the proof.
To do this, observe the following. 
\begin{itemize}
    \item By \Cref{prop:alternative_length},
    $l_{A/Q}(A/Q^i)$ is eventually polynomial with degree $d+1$.
    \item By \Cref{prop:alternative_length} again,
    $l_{A/\mm}(A/\mm^i)$ is eventually polynomial
    with degree $d'+1$.
    \item Since $A$ is Noetherian and $\sqrt Q=\mm$,
    we get that $Q\supseteq \mm^c$ for some $c$.
    This implies that there is a chain of two surjections
    $$A/\mm^{ci} \relbar\joinrel\twoheadrightarrow A/Q^i
    \relbar\joinrel\twoheadrightarrow A/\mm^i,$$
    which gives 
    $$l_{A/\mm}(A/\mm^{ci}) \geq l_{A/\mm}(A/Q^i) \geq 
    l_{A/\mm}(A/\mm^i).$$
    Thus, as $i\to\infty$,
    $l_{A/\mm}(A/Q^i)$, a polynomial of degree $d+1$,
    is squeezed between two polynomials of degree $d'+1$.
    Hence, $d+1=d'+1$, so $d=d'$.
\end{itemize}
Hence, the degree of $\dim_k \mm^{i+1}/\mm^i$
is $d\leq m-1$, so $\dimH A \leq m$.

\subsection{The Hilbert Dimension is at least the Krull Dimension}
\label{subsec:dimH_dimK}
In this section, we show that $\dimH A \geq \dimK A$.
The idea is to slice by a non-zero divisor $f\in\mm$
(i.e., replace $A$ with $A/(f)$).
We show that the slicing \textbf{must} decrease 
the Hilbert dimension by at least $1$,
but it decreases the Krull dimension by at most $1$.
This will allow us to use induction.
\begin{lemma}[Slicing decreases Hilbert dimension]
    \label{lem:hilbert_dim_slicing}
    Let $A$ be a Noetherian local ring with maximal ideal $\mm$.
    Let $f\in\mm$ that is not a zero divisor.
    Then,
    $$\dimH A/(f) \leq \dimH A - 1.$$
\end{lemma}
\begin{proof}
    Let $k=A/\mm$ and $A'=A/(f)$.
    Consider the short exact sequence
    $$0\longto \frac{(f)}{\mm^i\cap (f)} \longto \frac{A}{\mm^i} 
    \longto \frac{A'}{\mm^iA'}\longto 0,$$
    where the first map is embedding $(f)$ into $A$
    and the second map comes from quotient projection $A\to A'$.
    Taking dimensions gives
    \begin{align*}
    \dim_k\left(A'/\mm^i A'\right) 
    &= \dim_k(A/\mm^i) - \dim_k\left(
    \frac{(f)}{(f)\cap \mm^i}\right).
    \end{align*}
    By Artin-Rees Lemma (\Cref{thm:artin_rees}),
    the filtration $(f)\cap \mm^i$ is $\mm$-stable.
    This implies that there exists a constant $c$
    for which $(f)\cap \mm^i = f\mm^{i+c}$
    for all sufficiently large $i$.
    This implies $(f)/((f)\cap \mm^i) \simeq A/\mm^{i+c}$
    as $A$-module for all sufficiently large $i$.
    Therefore, 
    $$\dim_k(A'/\mm^i A') 
    = \dim_k(A/\mm^i) - \dim_k(A/\mm^{i+c}),$$
    which is indeed a polynomial of degree at most 
    $\dimH A-1$.
\end{proof}
We now proceed to prove the second part of the theorem.
\begin{proof}[Proof of $\dimH A\geq \dimK A$]
    We will do induction on $\dimH A$.

    If $\dimH A=0$, then $l(A/\mm^n)$ is constant for
    sufficiently large $n$.
    In particular, we have $\mm^n = \mm^{n+1} = \mm^{n+2}=\dots$,
    so by Nakayama (\Cref{cor:nakayama_local_ring}), 
    we get that $\mm^n=0$, so
    condition (d) of \Cref{thm:artinian_local_ring} applies.
    Hence, $A$ is Artinian, so $\dimK A=0$.

    Otherwise, assume $\dimH A>0$, then consider
    the chain of prime ideals
    $$\pp_0\subset\pp_1\subset\dots\subset\pp_s,$$
    and we want to show that $\dimH A\geq s$.

    The first thing we will do is to mod out by $\pp_0$.
    In particular, note that
    $$l(A/\mm^i) = l\left(\frac{A/\pp_0}{\mm^i}\right)
    \implies \dimH A \geq \dimH A/\pp_0.$$
    Therefore, by replacing $A$ with $A/\pp_0$,
    we can assume $\pp_0 = (0)$ and $A$ is a domain.

    The point of above is to make it easier
    to find a non-zero divisor for the lemma.
    In particular, for any $f\in\pp_1\setminus\{0\}$,
    $f$ is not a zero divisor.
    Now, we all we need is to cut the chain down to $A/(f)$.
    We set $\overline\pp_i = \pp_i/(f)$, so that
    $$\overline\pp_1\subset\overline\pp_2\subset\dots
    \subset \overline\pp_s \subset A/(f),$$
    so we get a chain of length $s-1$ in $A/f$.
    In particular, the induction hypothesis implies
    $\dimH A/(f)\geq s-1$. Hence,
    $$\dimH  A \geq \dimH A/(f) + 1 = s,$$
    so we are done.
\end{proof}
\subsection{The Krull Dimension is at least the Primary Ideal Dimension}
\label{subsec:dimK_dimP}
Finally, we prove the last inequality that
$\dimK A \geq \dimP A$.
In plain language, we have to find an $\mm$-primary
ideal generated by at most $\dimK A$ elements.
To do that, we use induction on the new measure of complexity of ideals, called \vocab{height}.
\begin{definition}[Height]
    For any ring $A$, the \vocab{height} 
    (also known as the \vocab{codimension}) of a prime ideal
    $\pp$, denoted $\height \pp$
    is the maximum length of a chain of primes 
    $\pp_0\subset\pp_1\subset\dots\subset\pp_n=\pp$
    ending at $\pp$.
    In other words, $\height\pp = \dim A_\pp$.

    The height of an ideal $I$ is defined as
    \notation{$\height I$}{height}
    $$\height I = \min_{\text{prime }\pp\supset I}
    \height\pp.$$
\end{definition}
We have the following easy properties of heights.
\begin{proposition}[Basic properties of height]
    \label{prop:height}
    If $A$ is a local ring with maximal ideal $\mm$,
    then we have the following.
    \begin{enumerate}[label=(\alph*)]
	\item $\height \pp + \dim A/\pp
	    \leq\dim A$.
	\item $\height I =0$
    if and only if $I$ is contained in a minimal prime of $A$.
	\item $\height I= \dim A$ if and only if
    $\sqrt I=\mm.$
    \end{enumerate}
\end{proposition}
\begin{proof}
    \begin{enumerate}[label=(\alph*)]
	\item By correspondence theorem,
    a chain in $A/\pp$ corresponds to a chain in $A$
    containing $\pp$, so we can stick the chain of $A/\pp$
    with the chain in $\height \pp$
    to get a chain of $A$, which must have 
    length at most $\dim A$.
    \item $\height I=0$ if and only if
    there is a prime $\pp$ such that $\height \pp=0$.
    This prime $\pp$ must be a minimal prime,
    so $I$ is contained in a minimal prime of $A$.
    \item We note that if $\height I=\dim A$,
    then for any $\pp\supset I$, we have the chain
    $\pp = \pp_{\dim A}\supset\pp_{\dim A-1}
    \supset\dots\supset \pp_0$,
    so $\pp$ is actually a maximal ideal.
    There is only one maximal ideal in $I$,
    so it implies that $\pp=\mm$, so $\sqrt I=\mm$.
    \qedhere
\end{enumerate}
\end{proof}
    
\begin{proof}[Proof of $\dimK A\geq \dimP A$]
    Let $d=\dimK A$.
   We need to find an $\mm$-primary ideal generated by 
   $\leq d$ elements.
   The key claim is the following:
   ``for each $i\leq d$, there exists
   $x_1,x_2,\dots,x_i\in\mm$
   such that $\height((x_1,\dots,x_i))\geq i$.''

   This key claim implies the theorem because
   we can find $x_1,\dots,x_{\dim A}$ satisfying
   the condition.
   Then, $Q=(x_1,\dots,x_{\dim A})$
   is $\mm$-primary by \Cref{prop:height} (c),
   implying that $\dimP A \leq d$.

   To prove the key claim,
   we find $x_1,\dots,x_i$ inductively.
   Assume that we have $x_1,\dots,x_{i-1}$ such that
   $\height(x_1,\dots,x_{i-1})\geq i-1$.
   Let $\pp_1,\dots,\pp_r$ be minimal primes over 
   $(x_1,\dots,x_{i-1})$ (i.e., minimal among ones
   containing $(x_1,\dots,x_{i-1})$); 
   there are finitely many because 
   \Cref{cor:finitely_minimal_primes} 
   applied on $A/(x_1,\dots,x_{i-1})$.
   These primes have height at most $i-1 \leq d-1$,
   so they are not equal to $\mm$.
   In particular, $\mm\neq \pp_j$ for all $j$,
   so by prime avoidance, \Cref{prob:prime_avoidance}, we have
   $\mm\neq\bigcup_{j=1}^r \pp_j$.
   
   Therefore, pick $x_i\in\mm$ such that $x_i\notin
   \bigcup\pp_j$ for all $j$.
   Then, for any prime $\qq\supset (x_1,\dots,x_{i-1},x_i)$,
   we have $\height(\qq)\geq i-1$,
   but $\qq\neq \pp_1,\dots,\pp_r$, so $\height(\qq)\geq i$.
\end{proof}
%
%
%
%
%
\section{Consequences of Dimension Theorem}
We now harvest some consequences of the dimension theorem.
\subsection{Dimension of Local Rings}
First, although dimension of Noetherian rings need not be finite,
we note the following clear corollary.
\begin{corollary}
    If $A$ is Noetherian local ring, then $\dim A$ is finite.
\end{corollary}
\begin{proof}
    Clear from either Hilbert dimension or primary
    ideal dimension.
\end{proof}
\begin{proposition}[Slicing by $f$]
    \label{prop:local_ring_slice}
    If $A$ is a Noetherian local ring with maximal ideal
    $\mm$ and $f\in\mm$
    is not a zero divisor, then 
    $$\dim A/(f) = \dim A - 1$$
\end{proposition}
\begin{proof}
    We showed the upper bound $\dim A/(f) \leq \dim A-1$
    over the course of proving dimension theorem 
    (\Cref{lem:hilbert_dim_slicing}).
    To prove the lower bound, we use primary ideal dimension:
    suppose that $(x_1,\dots,x_d)$ 
    generates an $\mm/(f)$-primary ideal.
    Then, its contraction to $A$, which is ideal
    $(f,x_1,\dots,x_d)$ would generate an $\mm$-primary ideal as well.
    Hence, $\dim A \leq \dim A/(f) + 1$, as desired.
\end{proof}
\begin{example}
    Suppose that you want to compute $\dim A$, where
    $$A = \frac{\CC[x,y,z]}{(xy+yz+zx)}$$
    using the above property. 
    We claim that for any maximal ideal $\mm$ of $A$,
    we have $\dim A_\mm = 2$. To see why,
    note that a maximal ideal $\mm$ of $A$
    is a contraction of a maximal ideal 
    $\frakn \subseteq \CC[x,y,z]$ such that $xy+yz+zx\in\frakn$.
    Note that $\CC[x,y,z]$ is a domain, so $xy+yz+zx$
    is not a zero divisor. Thus, by the above \Cref{prop:local_ring_slice},
    $$\dim A_\mm = \dim \frac{\CC[x,y,z]_{\frakn}}{(xy+yz+zx)}
    = \dim \CC[x,y,z]_\frakn - 1 = 3 - 1 = 2.$$
    Thus, taking supremum across all ideals $\mm$
    gives $\dim A = 2$.

    (This argument would have worked if $xy+yz+zx$
    is replaced with any $f\neq 0\in\CC[x,y,z]$.)
\end{example}
\begin{proposition}[Dimension of completion]
    If $A$ is a Noetherian local ring with maximal ideal $\mm$
    and $\widehat A$ is the $\mm$-adic completion,
    then $\dim \widehat  A = \dim A$.
\end{proposition}
\begin{proof}
    By \Cref{prop:completion_local}, $\widehat A$
    is local with maximal ideal $\widehat\mm$,
    which is a completion of $\mm$ as a module.
    From \Cref{cor:completion_isomorphism},
    we get that $\widehat \mm^i / \widehat \mm^{i+1}
    \simeq \mm^i / \mm^{i+1}$,
    so the Hilbert dimensions are the same.
\end{proof}
\subsection{Krull's Principal Ideal Theorem}
\begin{theorem}[Krull's Height Theorem]
\label{thm:krull_height}
Let $A$ be a Noetherian ring, and $I$ be an ideal with
$n$ generators. Then, $\height I \leq n$.
(Translation: every prime ideal minimal over $I$
has height at most $n$.)
\end{theorem}
\begin{proof}
    Let $\pp$ be a minimal prime among those containing $I$.
    Then, $\height(IA_\pp) = \dim A_\pp$, so by
    \Cref{prop:height}, $IA_\pp$ is $\pp A_\pp$-primary.
    Hence, using primary ideal dimension, we have
    $$\height \pp = \dim A_\pp  = \dimP A_\pp \leq n,$$
    so $\height I \leq n$.
\end{proof}
A special case of this is the following.
\begin{corollary}[Krull's Principal Ideal Theorem,
or Krull's Hauptidealsatz]
\label{cor:krull_principal}
Let $A$ be a Noetherian ring and $f\in A$. 
Let $\pp$ be a prime ideal minimal among those containing $f$.
\begin{alphenum}
    \item $\height\pp \leq 1$.
    \item If $f$ is not a zero divisor, then $\height \pp = 1$.
\end{alphenum}
\end{corollary}
\begin{proof}
    (a) follows immediately from \Cref{thm:krull_height}.
    To prove (b), note that if $\height\pp = 0$,
    then, $A_\pp$ has dimension $0$, hence Artinian
    (by \Cref{thm:artinian_local_ring} (b)).
    This means that $f\in \pp A_\pp$ is nilpotent
    (by \Cref{thm:artinian_local_ring} (d)), hence
    a zero divisor in $A$, a contradiction.
\end{proof}
\begin{remark}
    There are other proofs of Krull's principal ideal theorem
    and Krull's height theorem that does not need
    knowledge in this chapter.
    See for example, \cite[\S 12.3]{vakil}
\end{remark}
We also note the following useful property,
although the proof is deferred to \Cref{cor:dim_poly}.
\begin{proposition}[Dimension of polynomial ring]
    \label{prop:dim_poly}
    If $A$ is a Noetherian ring, then $\dim A[x] = \dim A + 1$.
\end{proposition}
%
%
%
%
%
%
\section{Regular Local Ring}
\subsection{Definition and Basic Properties}
For any Noetherian local ring $A$ with maximal ideal $\mm$
and $k=A/\mm$, we note that
\begin{align*}
    \dim A &\geq \#(\text{generators of }\mm) 
    \tag{primary ideal dimension} \\
    &\geq \dim_k(\mm/\mm^2)
    \tag{Nakayama; \Cref{cor:nakayama_generation}}.
\end{align*}
\begin{definition}[Regular local rings]
    A Noetherian local ring $(A,\mm)$ 
    is a \vocab{regular local ring} if the three above
    quantities are all equal:
    $$\dim A = \dim_{A/\mm} \mm/\mm^2
    = \#(\text{generators of }\mm).$$
    (Equivalently, the first two are equal.
    Also, we always assume that regular local rings
    are Noetherian.)
\end{definition}
\begin{example}[Regular local rings of dimension $0$ and $1$.]
\begin{itemize}
    \item $A$ is regular local with dimension $0$,
    if and only if $\mm=\mm^2$,
    which (by Nakayama; \Cref{cor:nakayama_local_ring})
    is equivalent to either $\mm=0$ or $A$ is a field.
    \item $A$ is a regular local ring with dimension 
    $1$ if and only if $A$ is a DVR 
    (by \Cref{thm:dvr} (d)).

    Thus, regular local ring generalizes the concept
    of DVR to higher-dimension.
\end{itemize}
\end{example}
\begin{example}[Regular local rings correspond to smooth objects]
Regular local rings correspond to rings of smooth varieties
localized at a point.
We will not explain this connection precisely.
Let us first confine ourself to a simpler example of hypersurface
$$A = \left(\frac{k[x_1,\dots,x_n]}{(f)}\right)_{(x_1,\dots,x_n)}$$
where $k$ is an algebraically closed field
and $f\in k[x_1,\dots,x_n]$ such that $f(0,\dots,0)=0$.
Thus, $\dim A = n-1$.
We write
$$f(x_1,\dots,x_n) = a_1x_1 + \dots + a_nx_n + 
(\text{higher degree terms}).$$

Let $\overline\mm$ be the maximal ideal of $A$,
which correspond to the maximal ideal $\mm$ 
of the localization 
$k[x_1,\dots,x_n]_{(x_1,\dots,x_n)}$. Then, we have
$$\frac{\overline\mm}{\overline\mm^2} 
= \frac{\mm}{(\mm^2, f)}
= \frac{\mm/\mm^2}{\Span(f\bmod \mm^2)}
\simeq \frac{\Span(e_1,\dots,e_n)}{\langle a_1e_1+\dots+a_ne_n=0
\rangle},$$
which has dimension $n-1$ if and only if 
at least one of $a_i$ is nonzero.

Equivalently, $A$ is a regular local ring if and only if
at least one partial derivative 
$\partial f / \partial x_i$ does not vanish at $(0,\dots,0)$.
\tcbline
In general, if $A$ is a localization at $(x_1,\dots,x_n)$ of
ring in form
$$k[x_1,\dots,x_n]/(f_1,\dots,f_r)$$
that has dimension $n-r$,
then $A$ is a regular local ring if and only if
the Jacobian matrix
$$\left(\left.\frac{\partial f_i}{\partial x_j}\right|_{(0,\dots,0)}\right)
_{1\leq i\leq r, 1\leq j\leq n}$$
has rank $n-r$. This is exactly the condition for
$k[x_1,\dots,x_n]/(f_1,\dots,f_r)$ to be a smooth variety
at $(0,\dots,0)$.

If $k$ is not algebraically closed, then 
maximal ideals may not be of the form $(x_1-a_1,\dots,x_n-a_n)$,
in which case checking regularity is harder
and beyond the scope of these notes.
Separability issues can arise of $k$ is not perfect.
\end{example}

Regular local rings are the nicest kinds of rings
that generalize to arbitrary dimension.
Unfortunately, studying regular local rings is quite
difficult. In the section, we will present something
that we can prove with all tools developed thus far.
The main highlight is that with mild conditions
(that holds in the case of any variety),
the $\mm$-adic completion of regular local ring $A$
must be isomorphic to $k[[x_1,\dots,x_d]]$.
Then, in \Cref{subsec:regular_local_hard_facts}, 
we state some nice facts that we \emph{cannot} prove.

First, we study the associated graded ring
$\gr_\mm A$ (\Cref{def:gr}), defined by
$$\gr_\mm A = \bigoplus_{i\geq 0} \mm^i/\mm^{i+1}$$
and show that they are the nicest as possible.
\begin{proposition}[$\gr$ of regular local ring]
    \label{prop:gr_regular_local}
    If $A$ is regular local with maximal ideal $\mm$
    and dimension $d$, then
    $\gr_\mm A \simeq k[x_1,\dots,x_d]$.

    \noindent Consequently, $\dim_{A/\mm} \mm^i/\mm^{i+1}
    = \binom{i+d-1}{d-1}$.
\end{proposition}
\begin{proof}
    Let $\mm = (t_1,\dots,t_d)$ be the minimal generating set.
    Then, we have a surjective map
    \begin{align*}
        k[x_1,\dots,x_d] &\to \gr_\mm A \\
        x_i &\mapsto t_i
    \end{align*}
    We have to show that this is injective.
    Suppose that a nonzero homogeneous $f$ of degree $s$
    maps to $0$.
    Then, the map collapse to a surjective map $k[x_1,\dots,x_d]/(f)
    \to \operatorname{gr}_\mm A$. Therefore,
    $$\dim \operatorname{gr}_\mm A
    \leq \dim \frac{k[x_1,\dots,x_d]}{(f)}
    = d-1.$$
    However, if $P = \bigoplus_{i\geq 1} \mm^i/\mm^{i+1}$
    be the ideal of positive degree,
    then $(\gr_\mm A)_P$ (localize at $P$)
    have the same associated graded ring as $A$,
    hence having the same Hilbert dimension as $A$.
    Thus, $\dim\gr_\mm A \geq \dim (\gr_\mm A)_P 
    = d$. This is a contradiction.
\end{proof}
\begin{corollary}
    Any regular local ring is a domain.
\end{corollary}
\begin{proof}
    Suppose $a,b\neq 0$.
    There exists $i,j$ such that
    $a\in\mm^i\setminus\mm^{i+1}$
    and $b\in\mm^j\setminus\mm^{j+1}$.
    Thus, $\overline a \in (\operatorname{gr}_\mm A)_i$
    and $\overline b\in (\operatorname{gr}_\mm A)_j$
    are both nonzero,
    so we have $\overline a\overline b \neq 0$.
    Thus, $ab\in\mm^{i+j}\setminus\mm^{i+j+1}$,
    so $ab\neq 0$.
\end{proof}
\begin{theorem}[Completion of regular local rings
is $k[[x_1,\dots,x_d]]$]
    \label{thm:completion_regular_local}
    Let $A$ be a regular local ring of dimension $d$,
    $\mm$ be the maximal ideal, and $k=A/\mm$.
    If every element $a\in k$ can be lifted to
    $\widetilde a\in A$ such that $a\equiv \widetilde a\pmod\mm$,
    then
    $$\widehat A \simeq k[[x_1,\dots,x_d]]$$
    (where $\widehat A$ denotes the $\mm$-adic completion).
\end{theorem}
Note that the lifting map condition holds whenever
$A = k[x_1,\dots,x_n]/I$ for some ideal $I$.
Thus, this theorem shows that completion of any smooth
variety is isomorphic to $k[[x_1,\dots,x_d]]$!
\begin{proof}
The idea is that from \Cref{prop:gr_regular_local},
we already have the isomorphism
$k[x_1,\dots,x_d]\to \gr_\mm A$. 
One needs to realize this as $\gr\psi$,
where $\psi : k[x_1,\dots,x_d] \to A$ is a map.
Then, we will be able to use \Cref{lem:gr_injective_surjective}
to conclude that $\widehat\psi : k[[x_1,\dots,x_d]] \to \widehat A$
is an isomorphism.

To do this, let $\mm = (t_1,\dots,t_d)$ as above. Define 
\begin{align*}
    \psi : k[x_1,\dots,x_d] &\to A \\
    a &\mapsto \widetilde a \quad\text{for all } a\in k \\
    x_i &\mapsto t_i
\end{align*}
Then, it is clear that $\gr\psi$ is the map defined
in \Cref{prop:gr_regular_local}, which is an isomorphism.
Thus, by \Cref{lem:gr_injective_surjective},
$\psi$ is an isomorphism.
\end{proof}
It's then natrual to ask whether all \textbf{complete} 
regular local rings are of the form $k[[x_1,\dots,x_n]]$.
(A ring is complete if it arises from a completion.)
It turns out that the answer is no;
there is a whole other family of regular local rings.
\begin{example}
    Let $B$ be a complete DVR (e.g., $B=\ZZ_p$).
    Let $\frakn=(\pi)$ be its maximal ideal and $B/\frakn=k$.
    We claim that
    $$A = B[[x_1,\dots,x_d]]$$
    is regular local of dimension $d+1$ with
    maximal ideal $\mm = (\pi,x_1,\dots,x_d)$. 
    This is because one can verify that
    $\gr_\mm A = k[\overline \pi, \overline x_1,
    \dots,\overline x_d].$
    For example, $\ZZ_p[[x,y]]$ is regular local
    which is not of the above form.
\end{example}

\subsection{Harder Facts about Regular Local Rings}
\label{subsec:regular_local_hard_facts}
Regular local rings enjoy further nice properties.
To conclude this section, we state some harder
facts about regular local rings that we will 
not be able to prove.

\begin{theorem}[Auslander–Buchsbaum]
    If $A$ is a regular local ring, then $A$ is a UFD
    (hence integrally closed).
\end{theorem}
See \cite[Thm. 19.19]{eisenbud} for the proof.
Note that the weaker fact that $A$ is integrally
closed is doable by elementary methods.

We also have a surprisingly difficult theorem by Serre.
\begin{theorem}[Serre]
    If $A$ is a regular local ring,
    then for any prime ideal $\pp$, the localization
    $A_\pp$ is still regular.
\end{theorem}
See \cite[Cor. 19.14]{eisenbud} for the proof.

These two theorems signify the next step of the subject,
which is to introduce \textbf{homological methods}.
Roughly speaking, homological methods introduces 
construction originally developed for algebraic topology
(e.g., chain complexes, homology, etc.).
These constructions allow us to cleanly manipulate
a sheer number of spaces and maps,
which then allow us to prove more difficult theorems.
You will see a glimpse of those homological techniques
in \Cref{chap:flatness}.

\section{Problems}
\begin{problem}
Let $A$ and $B$ be domains with $\dim A=\dim B$.
Prove that any surjective map $\phi : A\to B$ is an isomorphism.
\end{problem}
\begin{problem}
\label{prob:dim_pathology}
For a ring $A$, the property that ``$A$ is Noetherian''
and ``$\dim A$ is finite'' seem related in some ways.
However, neither one imply the other.
The following problems provide a counterexamples.
\begin{alphenum}
\item Consider $A=\CC[x_1,x_2,\dots]$. Let 
$$\pp_i = (x_{i^2}, x_{i^2+1}, x_{i^2+2}, \dots, 
x_{(i+1)^2+1})$$
and let $S=A\setminus \bigcup_{i=1}^\infty \pp_i$,
which is a multiplicative set.
Thus, $S^{-1}A$ is Noetherian 
(by \Cref{prop:noetherian_ring_properties}).
Prove that $\dim S^{-1}A = \infty$.

Hence $A$ is Noetherian does not imply that $\dim A$
is finite.
This counterexample was first discovered by Nagata.
\item Prove that the ring of algebraic integers 
has dimension $1$ (and we know that it is not Noetherian by 
\Cref{prob:algebraic_integer}).
Hence, $\dim A$ is finite does not imply that 
$A$ is Noetherian.
\end{alphenum}
\end{problem}
\begin{problem}
For each of the rings $A$ and maximal ideal $\mm$ in \Cref{prob:compute_gr}, compute the Hilbert dimension of $A_\mm$.
\end{problem}
\begin{problem}
Let $A=\CC[x,y]/(y^2-x^3-x)$ and $\mm=(x,y)$.
Find an $\mm$-primary ideal of $A_\mm$ that is principal (hence showing that $\dim A_\mm \leq 1$).
\end{problem}
\begin{problem}
Prove that a domain $A$ is a UFD if and only if 
every prime ideal with height $1$ is principal.
\end{problem}
\begin{problem}
Prove that the ring
$$\frac{\CC[w,x,y,z]}{(wz-xy, wy-x^2, xz-y^2)}$$
has dimension $2$ (despite it is cut out by $3$ equations).
\end{problem}
\begin{problem}
Prove that 
$$\left(\frac{\CC[x,y,z]}{(x^2-yz)}\right)_{(x,y,z)}$$
is not a regular local ring, even though \Cref{prob:integrally_closed} shows that it is integrally closed.
\end{problem}

The next two problems concern integrally closed domains.
\begin{problem}[Serre's R1 + S2 Criterion]
Let $A$ be a domain. Prove that $A$ is integrally closed
if and only if both of the following holds.
\begin{enumerate}[label=(\roman*)]
\item $A_\pp$ is a DVR for all prime ideal $\pp$ with height $1$.
\item For any $f\in A$, all associated primes of 
$A$-module $A/(f)$ are minimal over $(f)$
(and thus have codimension $1$).
\end{enumerate}

(Hint: to prove that integrally closed implies (ii),
reduce to the local case and mimic the proof of \Cref{thm:dvr} (b) $\Rightarrow$ (c).
To prove that (i) and (ii) implies integrally closed,
use \Cref{prob:ass_denominator}.)
\end{problem}
\begin{problem}[Algebraic Hartogs Lemma]
Let $A$ be an integrally closed domain. Prove that
$$A = \bigcap_{\pp \text{ height }1} A_\pp.$$
\end{problem}

\begin{problem}[Cohen-Macaulay Rings]
\label{prob:cohen_macaulay}
Cohen-Macaulay is a condition for local rings to 
be ``nice'' in a certain way that excludes some pathology,
but includes regular local rings and many non-smooth examples.
This problem gives an introduction to 
the theory of Cohen-Macaulay rings.

Let $A$ be a Noetherian local ring. 
A \vocab{regular sequence} is a sequence $f_1,\dots,f_n\in A$
such that $(f_1,\dots,f_n)\neq A$ 
and $x_i$ is not a zero divisor in 
$A/(f_1,\dots,f_{i-1})$. 
You should think of $f_1,\dots,f_n$ as equations, 
and we are slicing by the equations $f_i$. 

The \vocab{depth} of $A$, denoted $\operatorname{depth} A$ 
is the maximum length of regular sequences of $A$. 
A ring $A$ is \vocab{Cohen-Macaulay} if 
$\operatorname{depth} A = \dim A$.
Informally, this means that you can slice $\dim A$ equations.
\begin{alphenum}
\item Prove that $\operatorname{depth} A \leq \dim A$ 
for every ring $A$.
\item Suppose that $A$ has dimension $1$. 
Prove that $A$ is Cohen-Macaulay 
if and only all associated primes of $A$ 
are minimal.
(In particular, by \Cref{prob:ass_reduced},
any reduced Noetherian local ring with dimension $1$ 
is Cohen-Macaulay.) 
Given an example of a ring that is not Cohen-Macaulay.
\item Prove that regular local rings are Cohen-Macaulay.
\item Prove that the local ring
$$\left(\frac{\CC[x,y,z]}{(x^2-yz)}\right)_{(x,y,z)}$$
is Cohen-Macaulay.
\end{alphenum}
\end{problem}